\begin{apartats}

\apartat[1,25]{1}
Aplicant la definició de derivada, calcula la funció derivada de
\[
f(x)=x^3 .
\]

\begin{solucio}
Com que $(x+h)^3=x^3+3x^2h+3xh^2+h^3$,
\[
f'(x)=\lim_{h\to0}\frac{(x+h)^3-x^3}{h}
=\lim_{h\to0}\frac{3x^2h+3xh^2+h^3}{h}
=\lim_{h\to0}\left(3x^2+3xh+h^2\right)=3x^2 .
\]
\end{solucio}

\begin{nomesllarg}
\apartat{0,75}
Aplicant la definició de derivada, calcula la funció derivada de
\[
g(x)=-\frac{3}{x} .
\]

\begin{solucio}
\[
g'(x)=\lim_{h\to0}\frac{-\frac{3}{x+h}+\frac{3}{x}}{h}
=\lim_{h\to0}\frac{-3x+3(x+h)}{h\,x(x+h)}
=\lim_{h\to0}\frac{3}{x(x+h)}=\frac{3}{x^2} ,
\]
per a $x\neq0$.
\end{solucio}
\end{nomesllarg}

\apartat[1,25]{0,75}
Aplicant la definició de derivada, calcula la funció derivada de
\[
k(x)=2\sqrt{x} .
\]

\begin{solucio}
Amb el conjugat:
\[
k'(x)=\lim_{h\to0}\frac{2\sqrt{x+h}-2\sqrt{x}}{h}
=\lim_{h\to0}\frac{2\left((x+h)-x\right)}{h\left(\sqrt{x+h}+\sqrt{x}\right)}
=\lim_{h\to0}\frac{2}{\sqrt{x+h}+\sqrt{x}}=\frac{1}{\sqrt{x}} ,
\]
per a $x>0$.
\end{solucio}

\end{apartats}
